\documentclass[10pt,a4paper]{amsart}
\usepackage{amsmath,amssymb,amsthm,graphicx}
\usepackage{verbatim}
\newenvironment{reference}{\comment}{\endcomment}
%\newenvironment{reference}{}{}
\newenvironment{slogan}{\comment}{\endcomment}
\newenvironment{history}{\comment}{\endcomment}

% For commutative diagrams we use Xy-pic
\usepackage[all]{xy}

% We use 2cell for 2-commutative diagrams.
\xyoption{2cell}
% Theorem environments.
%
\theoremstyle{plain}
\newtheorem{theorem}[subsection]{Theorem}
\newtheorem{proposition}[subsection]{Proposition}
\newtheorem{lemma}[subsection]{Lemma}

\theoremstyle{definition}
\newtheorem{definition}[subsection]{Definition}
\newtheorem{example}[subsection]{Example}
\newtheorem{exercise}[subsection]{Exercise}
\newtheorem{situation}[subsection]{Situation}

\theoremstyle{remark}
\newtheorem{remark}[subsection]{Remark}
\newtheorem{remarks}[subsection]{Remarks}

\numberwithin{equation}{subsection}

% Macros
%
\def\lim{\mathop{\mathrm{lim}}\nolimits}
\def\colim{\mathop{\mathrm{colim}}\nolimits}
\def\Spec{\mathop{\mathrm{Spec}}}
\def\Hom{\mathop{\mathrm{Hom}}\nolimits}
\def\Ext{\mathop{\mathrm{Ext}}\nolimits}
\def\SheafHom{\mathop{\mathcal{H}\!\mathit{om}}\nolimits}
\def\SheafExt{\mathop{\mathcal{E}\!\mathit{xt}}\nolimits}
\def\Sch{\mathit{Sch}}
\def\Mor{\mathop{\mathrm{Mor}}\nolimits}
\def\Ob{\mathop{\mathrm{Ob}}\nolimits}
\def\Sh{\mathop{\mathit{Sh}}\nolimits}
\def\NL{\mathop{N\!L}\nolimits}
\def\CH{\mathop{\mathrm{CH}}\nolimits}
\def\proetale{{pro\text{-}\acute{e}tale}}
\def\etale{{\acute{e}tale}}
\def\QCoh{\mathit{QCoh}}
\def\Ker{\mathop{\mathrm{Ker}}}
\def\Im{\mathop{\mathrm{Im}}}
\def\Coker{\mathop{\mathrm{Coker}}}
\def\Coim{\mathop{\mathrm{Coim}}}

% Boxtimes
%
\DeclareMathSymbol{\boxtimes}{\mathbin}{AMSa}{"02}

%
% Macros for moduli stacks/spaces
%
\def\QCohstack{\mathcal{QC}\!\mathit{oh}}
\def\Cohstack{\mathcal{C}\!\mathit{oh}}
\def\Spacesstack{\mathcal{S}\!\mathit{paces}}
\def\Quotfunctor{\mathrm{Quot}}
\def\Hilbfunctor{\mathrm{Hilb}}
\def\Curvesstack{\mathcal{C}\!\mathit{urves}}
\def\Polarizedstack{\mathcal{P}\!\mathit{olarized}}
\def\Complexesstack{\mathcal{C}\!\mathit{omplexes}}
% \Pic is the operator that assigns to X its picard group, usage \Pic(X)
% \Picardstack_{X/B} denotes the Picard stack of X over B
% \Picardfunctor_{X/B} denotes the Picard functor of X over B
\def\Pic{\mathop{\mathrm{Pic}}\nolimits}
\def\Picardstack{\mathcal{P}\!\mathit{ic}}
\def\Picardfunctor{\mathrm{Pic}}
\def\Deformationcategory{\mathcal{D}\!\mathit{ef}}
\usepackage{cleveref}
\crefname{originalsection}{original section}{original sections}
\Crefname{originalsection}{Original section}{Original sections}
\makeatletter
\newlabel{homology-example-graded-vector-spaces}{{0FFX}{source}{}{sourcepointer-1}{}}
\newlabel{homology-example-graded-vector-spaces@cref}{{[originalsection][0][]0FFX}{[1][0][]source}{}{}{}}
\makeatother
\crefname{originallemma}{original lemma}{original lemmas}
\Crefname{originallemma}{Original lemma}{Original lemmas}
\makeatletter
\newlabel{homology-lemma-left-dual-graded-vector-spaces}{{0FFV}{source}{}{sourcepointer-2}{}}
\newlabel{homology-lemma-left-dual-graded-vector-spaces@cref}{{[originallemma][0][]0FFV}{[1][0][]source}{}{}{}}
\makeatother
\crefname{originalreference}{original reference}{original references}
\Crefname{originalreference}{Original reference}{Original references}
\makeatletter
\newlabel{lemma-composition-correspondences}{{0FG0}{source}{}{sourcepointer-3}{}}
\newlabel{lemma-composition-correspondences@cref}{{[originalreference][0][]0FG0}{[1][0][]source}{}{}{}}
\makeatother
\crefname{originalreference}{original reference}{original references}
\Crefname{originalreference}{Original reference}{Original references}
\makeatletter
\newlabel{lemma-dual}{{08I1}{source}{}{sourcepointer-4}{}}
\newlabel{lemma-dual@cref}{{[originalreference][0][]08I1}{[1][0][]source}{}{}{}}
\makeatother
\crefname{originalreference}{original reference}{original references}
\Crefname{originalreference}{Original reference}{Original references}
\makeatletter
\newlabel{lemma-functor-and-cycles}{{0FG5}{source}{}{sourcepointer-5}{}}
\newlabel{lemma-functor-and-cycles@cref}{{[originalreference][0][]0FG5}{[1][0][]source}{}{}{}}
\makeatother
\crefname{originalreference}{original reference}{original references}
\Crefname{originalreference}{Original reference}{Original references}
\makeatletter
\newlabel{lemma-prep-dual}{{0FG8}{source}{}{sourcepointer-6}{}}
\newlabel{lemma-prep-dual@cref}{{[originalreference][0][]0FG8}{[1][0][]source}{}{}{}}
\makeatother
\crefname{originalreference}{original reference}{original references}
\Crefname{originalreference}{Original reference}{Original references}
\makeatletter
\newlabel{remark-replace-cup-product-classical}{{0FGU}{source}{}{sourcepointer-7}{}}
\newlabel{remark-replace-cup-product-classical@cref}{{[originalreference][0][]0FGU}{[1][0][]source}{}{}{}}
\makeatother
\title{Producing the classical axioms of a Weil cohomology theory from a monoidal functor on motives}
\author{The Stacks Project Authors}
\date{stacks-project:weil.tex:local-snapshot-20261012; GNU FDL 1.2 or later}
\begin{document}
\maketitle
\noindent\textit{Editorial note.} All delivered mathematical text is the author's own, copied from the pinned Stacks Project chapter file weil.tex; no statement or proof step was rewritten, reordered or removed.\par\medskip
\noindent\textit{Editorial note.} The counted proof is the author's complete argument: the construction of the graded commutative cohomology algebra, the identification of the twists of the coefficients through chosen isomorphisms, the cycle class and degree maps, and the verification of the duality, product and functoriality axioms.\par\medskip
\noindent\textit{Editorial note.} The label lemma-from-functor-to-weil-classical names the counted statement, and the proof environment immediately following it in the source is delivered verbatim as its proof.\par\medskip
\noindent\textit{Editorial note.} The counted statement is carried by the author's own lemma environment in the chapter source.\par\medskip
\noindent\textit{Editorial note.} Pre-open gate: python operations/agent-authoring/tools/precard\_gate.py --source weil.tex --label lemma-from-functor-to-weil-classical --start 1857 --end 2027 --local-id 28135 returned RESULT=OK (L28135 free, occupancy FREE, public interval clear). Pre-assembly tag self-check: canonical tag resolved and not present in reference/public-items.jsonl.\par\medskip
\noindent\textit{Editorial note.} The statement and proof are copied from the pinned chapter file; the cited inputs are kept as external pointers with their canonical Stacks tags.\par\medskip
\noindent\textit{Allowed standard prerequisites (editorial).}\begin{itemize}
\item Correspondences as morphism groups between motives over an algebraically closed field. Same chapter, the lemma identifying correspondence groups with morphism groups of motives and with rational Chow groups. Provides the morphism groups whose images under the functor give the cycle class map and the degree map.
\item Duality of the twisted motive and its evaluation as the class of the diagonal. Same chapter, the lemma exhibiting the twisted motive of a smooth projective variety as a left dual, with the evaluation given by the class of the diagonal. Gives the nondegenerate pairings on cohomology and their identification with cup product followed by the degree map.
\item The projection formula through the duality preparation lemma. Same chapter, the preparation lemma for the duality diagram, together with the lemmas relating pushforwards of cycles to compositions of correspondences. Verifies the compatibility of the cycle class map with pushforwards and the projection formula along a morphism.
\end{itemize}\medskip
\ifdefined\hypertarget\hypertarget{sourcepointer-1}{}\fi
\noindent\textit{Source reference.} Original section 0FFX: sources/stacks-chow/tags.txt, source lines 19302--19302. Prerequisite pointer invoked at proof line(s) 1889; it names a standard fact of a sibling chapter that the counted argument uses, and every step of the counted proof built on it is retained verbatim in the delivered body.\par\medskip
\ifdefined\hypertarget\hypertarget{sourcepointer-2}{}\fi
\noindent\textit{Source reference.} Original lemma 0FFV: sources/stacks-chow/tags.txt, source lines 19300--19300. Prerequisite pointer invoked at proof line(s) 1940; it names a standard fact of a sibling chapter that the counted argument uses, and every step of the counted proof built on it is retained verbatim in the delivered body.\par\medskip
\ifdefined\hypertarget\hypertarget{sourcepointer-3}{}\fi
\noindent\textit{Source reference.} Original reference 0FG0: sources/stacks-weil/pinned-original-source.tex, source lines 441--441. Prerequisite pointer invoked at proof line(s) 1906; it names a standard fact of the same chapter that the counted argument uses, and every step of the counted proof built on it is retained verbatim in the delivered body.\par\medskip
\ifdefined\hypertarget\hypertarget{sourcepointer-4}{}\fi
\noindent\textit{Source reference.} Original reference 08I1: sources/stacks-weil/pinned-original-source.tex, source lines 1186--1186. Prerequisite pointer invoked at proof line(s) 1936; it names a standard fact of the same chapter that the counted argument uses, and every step of the counted proof built on it is retained verbatim in the delivered body.\par\medskip
\ifdefined\hypertarget\hypertarget{sourcepointer-5}{}\fi
\noindent\textit{Source reference.} Original reference 0FG5: sources/stacks-weil/pinned-original-source.tex, source lines 661--661. Prerequisite pointer invoked at proof line(s) 1966; it names a standard fact of the same chapter that the counted argument uses, and every step of the counted proof built on it is retained verbatim in the delivered body.\par\medskip
\ifdefined\hypertarget\hypertarget{sourcepointer-6}{}\fi
\noindent\textit{Source reference.} Original reference 0FG8: sources/stacks-weil/pinned-original-source.tex, source lines 777--777. Prerequisite pointer invoked at proof line(s) 1980; it names a standard fact of the same chapter that the counted argument uses, and every step of the counted proof built on it is retained verbatim in the delivered body.\par\medskip
\ifdefined\hypertarget\hypertarget{sourcepointer-7}{}\fi
\noindent\textit{Source reference.} Original reference 0FGU: sources/stacks-weil/pinned-original-source.tex, source lines 1653--1653. Prerequisite pointer invoked at proof line(s) 2013; it names a standard fact of the same chapter that the counted argument uses, and every step of the counted proof built on it is retained verbatim in the delivered body.\par\medskip
\begin{lemma}
\label{lemma-from-functor-to-weil-classical}
Let $k$ be an algebraically closed field. Let $F$ be a field of
characteristic $0$. Consider a $\mathbf{Q}$-linear functor
$$
G : M_k \longrightarrow \text{graded }F\text{-vector spaces}
$$
of symmetric monoidal categories such that $G(\mathbf{1}(1))$
is nonzero only in degree $-2$. Then we obtain data (D1), (D2), (D3)
satisfying all of (A), (B), (C) except for possibly (A)(c) and (A)(d).
\end{lemma}

\begin{proof}
We obtain a contravariant functor from the category of smooth
projective varieties to the category of graded $F$-vector spaces
by setting $H^*(X) = G(h(X))$. By assumption we have a canonical
isomorphism
$$
H^*(X \times Y) = G(h(X \times Y)) = G(h(X) \otimes h(Y)) =
G(h(X)) \otimes G(h(Y)) = H^*(X) \otimes H^*(Y)
$$
compatible with pullbacks. Using pullback along the diagonal
$\Delta : X \to X \times X$ we obtain a canonical map
$$
H^*(X) \otimes H^*(X) = H^*(X \times X) \to H^*(X)
$$
of graded vector spaces compatible with pullbacks.
This defines a functorial graded $F$-algebra structure on
$H^*(X)$. Since $\Delta$ commutes with the commutativity
constraint $h(X) \otimes h(X) \to h(X) \otimes h(X)$ (switching the factors)
and since $G$ is a functor of symmetric monoidal categories (so compatible with
commutativity constraints), and by our convention in
Homology, Example \ref{homology-example-graded-vector-spaces}
we conclude that $H^*(X)$ is a graded
commutative algebra! Hence we get our datum (D1).

\medskip\noindent
Since $\mathbf{1}(1)$ is invertible in the category of motives
we see that $G(\mathbf{1}(1))$ is invertible in the category of
graded $F$-vector spaces. Thus $\sum_i \dim_F G^i(\mathbf{1}(1)) = 1$.
By assumption we only get something nonzero in degree $-2$ and we may
choose an isomorphism $F[2] \to G(\mathbf{1}(1))$ of graded $F$-vector spaces.
Here and below $F[n]$ means the graded $F$-vector space which has
$F$ in degree $-n$ and zero elsewhere. Using compatibility with
tensor products, we find for all $n \in \mathbf{Z}$ an isomorphism
$F[2n] \to G(\mathbf{1}(n))$ compatible with tensor products.

\medskip\noindent
Let $X$ be a smooth projective variety. By
Lemma \ref{lemma-composition-correspondences} we have
$$
\CH^r(X) \otimes \mathbf{Q} = \text{Corr}^r(\Spec(k), X) =
\Hom(\mathbf{1}(-r), h(X))
$$
Applying the functor $G$ we obtain
$$
\gamma :
\CH^r(X) \otimes \mathbf{Q} \longrightarrow
\Hom(G(\mathbf{1}(-r)), H^*(X)) = H^{2r}(X)
$$
This is the datum (D2).

\medskip\noindent
Let $X$ be a smooth projective variety of dimension $d$. By
Lemma \ref{lemma-composition-correspondences} we have
$$
\Mor(h(X)(d), \mathbf{1}) = \Mor((X, 1, d), (\Spec(k), 1, 0)) =
\text{Corr}^{-d}(X, \Spec(k)) = \CH_d(X)
$$
Thus the class of the cycle $[X]$ in $\CH_d(X)$ defines a morphism
$h(X)(d) \to \mathbf{1}$. Applying $G$ we obtain
$$
H^*(X) \otimes F[-2d] = G(h(X)(d)) \longrightarrow G(\mathbf{1}) = F
$$
This map is zero except in degree $0$ where we obtain
$\int_X : H^{2d}(X) \to F$. This is the datum (D3).

\medskip\noindent
Let $X$ be a smooth projective variety of dimension $d$.
By Lemma \ref{lemma-dual}
we know that $h(X)(d)$ is a left dual to $h(X)$. Hence
$G(h(X)(d)) = H^*(X) \otimes F[-2d]$ is a left dual to
$H^*(X)$ in the category of graded $F$-vector spaces.
By Homology, Lemma \ref{homology-lemma-left-dual-graded-vector-spaces}
we find that $\sum_i \dim_F H^i(X) < \infty$ and that
$\epsilon : h(X)(d) \otimes h(X) \to \mathbf{1}$ produces
nondegenerate pairings $H^{2d - i}(X) \otimes_F H^i(X) \to F$.
In the proof of Lemma \ref{lemma-dual} we have seen that
$\epsilon$ is given by $[\Delta]$ via the identifications
$$
\Hom(h(X)(d) \otimes h(X), \mathbf{1}) =
\text{Corr}^{-d}(X \times X, \Spec(k)) =
\CH_d(X \times X)
$$
Thus $\epsilon$ is the composition of $[X] : h(X)(d) \to \mathbf{1}$
and $h(\Delta)(d) : h(X)(d) \otimes h(X) \to h(X)(d)$. It follows
that the pairings above are given by cup product followed by
$\int_X$. This proves axiom (A) parts (a) and (b).

\medskip\noindent
Axiom (B) follows from the assumption that $G$ is compatible
with tensor structures and our construction of the cup product above.

\medskip\noindent
Axiom (C). Our construction of $\gamma$ takes a cycle $\alpha$ on $X$,
interprets it as a correspondence $a$ from $\Spec(k)$ to $X$ of some degree,
and then applies $G$. If $f : Y \to X$ is a morphism of smooth projective
varieties, then $f^!\alpha$ is the pushforward (!) of $\alpha$
by the correspondence $[\Gamma_f]$ from $X$ to $Y$, see
Lemma \ref{lemma-functor-and-cycles}. Hence
$f^!\alpha$ viewed as a correspondence from $\Spec(k)$ to $Y$
is equal to $a \circ [\Gamma_f]$, see
Lemma \ref{lemma-composition-correspondences}.
Since $G$ is a functor, we conclude
$\gamma$ is compatible with pullbacks, i.e., axiom (C)(a) holds.

\medskip\noindent
Let $f : Y \to X$ be a morphism of smooth projective varieties and
let $\beta \in \CH^r(Y)$ be a cycle on $Y$. We have to show that
$$
\int_Y \gamma(\beta) \cup f^*c = \int_X \gamma(f_*\beta) \cup c
$$
for all $c \in H^*(X)$. Let $a, a^t, \eta_X, \eta_Y, [X], [Y]$
be as in Lemma \ref{lemma-prep-dual}.
Let $b$ be $\beta$ viewed as a correspondence from $\Spec(k)$ to $Y$
of degree $r$. Then $f_*\beta$ viewed as a correspondence from
$\Spec(k)$ to $X$ is equal to $a^t \circ b$, see
Lemmas \ref{lemma-functor-and-cycles} and
\ref{lemma-composition-correspondences}.
The displayed equality above holds if we can show that
$$
h(X) = \mathbf{1} \otimes h(X)
\xrightarrow{b \otimes 1}
h(Y)(r) \otimes h(X)
\xrightarrow{1 \otimes a}
h(Y)(r) \otimes h(Y)
\xrightarrow{\eta_Y}
h(Y)(r)
\xrightarrow{[Y]}
\mathbf{1}(r - e)
$$
is equal to
$$
h(X) = \mathbf{1} \otimes h(X)
\xrightarrow{a^t \circ b \otimes 1}
h(X)(r + d - e) \otimes h(X)
\xrightarrow{\eta_X}
h(X)(r + d - e)
\xrightarrow{[X]}
\mathbf{1}(r - e)
$$
This follows immediately from Lemma \ref{lemma-prep-dual}.
Thus we have axiom (C)(b).

\medskip\noindent
To prove axiom (C)(c) we use the discussion in
Remark \ref{remark-replace-cup-product-classical}.
Hence it suffices to prove that $\gamma$ is compatible with
exterior products. Let $X$, $Y$ be smooth projective varieties and
let $\alpha$, $\beta$ be cycles on them. Denote
$a$, $b$ the corresponding correspondences from $\Spec(k)$ to
$X$, $Y$. Then $\alpha \times \beta$ corresponds to the
correspondence $a \otimes b$ from $\Spec(k)$ to $X \otimes Y = X \times Y$.
Hence the requirement follows from the fact that $G$ is
compatible with the tensor structures on both sides.

\medskip\noindent
Axiom (C)(d) follows because the cycle $[\Spec(k)]$
corresponds to the identity morphism on $h(\Spec(k))$.
This finishes the proof of the lemma.
\end{proof}

\end{document}
